\subchapter
{Basic Buildroot usage}
{Objectives:
  \begin{itemize}
  \item Get Buildroot
  \item Configure a minimal system with Buildroot for the BeaglePlay
  \item Do the build
  \item Prepare the BeaglePlay for usage
  \item Flash and test the generated system
  \end{itemize}
}

\section{Setup}

Go to the \code{$HOME/__SESSION_NAME__-labs/} directory.

As specified in the Buildroot
manual\footnote{\url{https://buildroot.org/downloads/manual/manual.html\#requirement-mandatory}},
Buildroot requires a few packages to be installed on your
machine. Let's install them using Ubuntu's package manager:

\begin{bashinput}
sudo apt install sed make binutils gcc g++ bash patch \
  gzip bzip2 perl tar cpio python3 unzip rsync wget libncurses-dev
\end{bashinput}

\section{Download Buildroot}

Since we're going to do Buildroot development, let's clone the
Buildroot source code from its Git repository:

\begin{bashinput}
git clone https://gitlab.com/buildroot.org/buildroot.git
\end{bashinput}

Go into the newly created \code{buildroot} directory.

We're going to start a branch from the {\em 2025.02.6} Buildroot
release, with which this training has been tested.

\begin{bashinput}
git checkout -b bootlin 2025.02.6
\end{bashinput}

\section{Configuring Buildroot}

If you look under \code{configs/}, you will see that there are
ready-to-use Buildroot configuration files for various platforms.
However, since we want to learn about Buildroot, we'll start our own
configuration from scratch!

Start the Buildroot configuration utility:

\begin{bashinput}
make menuconfig
\end{bashinput}

Of course, you're free to try out the other configuration utilities
\code{nconfig}, \code{xconfig} or \code{gconfig}.

Now, let's do the configuration:

\begin{itemize}

\item \code{Target Options} menu

  \begin{itemize}

  \item The BeaglePlay uses a Texas Instruments AM625x SoC, which
    contains four ARM Cortex-A53 cores. This is a 64-bit ARM platform,
    so select \code{AArch64 (little endian)} as the target architecture.

  \item Select \code{cortex-A53} as the \code{Target Architecture
      Variant}.

  \item The other parameters can be left to their default values.

  \end{itemize}

\item We don't have anything special to change in the
  \code{Build options} menu, but take nonetheless this opportunity to
  visit this menu, and look at the available options. Each option has
  a help text that tells you more about the option.

\item \code{Toolchain} menu

  \begin{itemize}

  \item By default, Buildroot builds its own toolchain. This takes
    quite a bit of time, and for \code{AArch64} platforms, there is a
    pre-built toolchain provided by Bootlin. We'll use it through the
    {\em external toolchain} mechanism of Buildroot. Select
    \code{External toolchain} as the \code{Toolchain type}. Do not
    hesitate however to look at the available options when you select
    \code{Buildroot toolchain} as the \code{Toolchain type}.

  \item Select \code{Bootlin toolchains} as the \code{Toolchain}. It
    will automatically select the \code{aarch64 glibc bleeding-edge
      2024.05-1} variant, which is fine for our needs. Buildroot can
    either use pre-defined toolchains such as the ones provided by ARM
    or Bootlin, or custom toolchains (either downloaded from a given
    location, or pre-installed on your machine).

  \end{itemize}

\item \code{System configuration} menu

  \begin{itemize}

  \item For our basic system, we don't need a lot of custom {\em
      system configuration} for the moment. So take some time to look
    at the available options, and put some custom values for the
    \code{System hostname}, \code{System banner} and \code{Root
      password}.

  \end{itemize}

\item \code{Kernel} menu

  \begin{itemize}

  \item We obviously need a Linux kernel to run on our platform, so
    enable the \code{Linux kernel} option.

  \item By default, the most recent Linux kernel version available at
    the time of the Buildroot release is used. In our case, we want to
    use a specific version, to make sure our build is reproducible. So
    select \code{Custom version} as the \code{Kernel version}, and
    enter \code{6.12.47} in the \code{Kernel version} text field that
    appears.

  \item Now, we need to define which kernel configuration to
    use. We'll start by using a default configuration provided within
    the kernel sources themselves. Select \code{Use the architecture
      default configuration} as the \code{Kernel configuration}. On
    ARM64, this uses the unified \code{defconfig} which supports all
    ARM64 platforms, including the TI AM62x.

  \item The \code{Kernel binary format} is the next option. On AArch64
    platforms, the kernel image format is \code{Image}. Make sure this
    option is selected.

  \item On ARM64, all platforms use the {\em Device Tree} to
    describe the hardware. The BeaglePlay is in this situation, so
    you'll have to enable the \code{Build a Device Tree Blob (DTB)}
    option. At
    \url{https://git.kernel.org/cgit/linux/kernel/git/torvalds/linux.git/tree/arch/arm64/boot/dts/ti/?id=v6.12},
    you can see the list of Device Tree files for TI ARM64 platforms.
    The one for the BeaglePlay is \code{k3-am625-beagleplay.dts}.
    Back in Buildroot, enable \code{Build a Device Tree Blob (DTB)}
    and type \code{ti/k3-am625-beagleplay} as the \code{In-tree Device
    Tree Source file names}.

  \item The kernel configuration for this platform requires having
    OpenSSL available on the host machine. To avoid depending on the
    OpenSSL development files installed by your host machine Linux
    distribution, Buildroot can build its own version: just enable the
    \code{Needs host OpenSSL} option.

  \end{itemize}

\item \code{Target packages} menu. This is probably the most important
  menu, as this is the one where you can select amongst the 3000+
  available Buildroot packages which ones should be built and
  installed in your system. For our basic system, enabling
  \code{BusyBox} is sufficient and is already enabled by default, but
  feel free to explore the available packages. We'll have the
  opportunity to enable some more packages in the next labs.

\item \code{Filesystem images} menu. For now, keep only the \code{tar
    the root filesystem} option enabled. We'll take care separately of
  flashing the root filesystem on the SD card.

\item \code{Bootloaders} menu.

  \begin{itemize}

  \item We'll use the most popular ARM bootloader, {\em U-Boot}, so
    enable it in the configuration.

  \item Select \code{Kconfig} as the \code{Build system}. U-Boot is
    transitioning from a situation where all the hardware platforms
    were described in C header files to a system where U-Boot re-uses
    the Linux kernel configuration logic. Since we are going to use a
    recent enough U-Boot version, we are going to use the latter,
    called {\em Kconfig}.

  \item Use the custom version of U-Boot \code{2024.04}.

  \item The BeaglePlay uses the TI AM62x SoC from the K3 family. The
    U-Boot configuration for the A53 cores is
    \code{am62x_beagleplay_a53_defconfig}. So enter
    \code{am62x_beagleplay_a53} as the \code{Board defconfig}.

  \item Select \code{u-boot.img} as the \code{U-Boot binary format}.
    We only need Buildroot to produce \code{u-boot.img} (the second
    stage bootloader). The first stage bootloaders are provided
    pre-built (see note below).

  \item In \code{Custom make options}, enter
    \code{BINMAN_ALLOW_MISSING=1}. This allows the U-Boot build to
    complete even though we are not building the first-stage boot
    images (which require additional firmware blobs).

  \end{itemize}

\item {\bf Note:} The AM62x SoC has a complex multi-stage boot
    process involving multiple processor cores. Before the Cortex-A53
    cores can run U-Boot, a Cortex-R5F core must first initialize the
    SoC (configure the DDR memory, load TI Foundational Security
    firmware, etc.). The first-stage boot chain consists of:

    \begin{itemize}
    \item \code{tiboot3.bin}: the R5 U-Boot SPL combined with TI
      system firmware, which initializes DDR and starts the A53 cores.
    \item \code{tispl.bin}: a FIT image containing the A53 U-Boot SPL,
      ARM Trusted Firmware (TF-A), OP-TEE, and the TI Device Manager
      firmware. This is loaded by \code{tiboot3.bin} and handles the
      secure boot chain before loading \code{u-boot.img}.
    \end{itemize}

    Building these first-stage images requires TI's downstream U-Boot
    fork, a separate 32-bit ARM toolchain (\code{arm-none-eabi-}), and
    TI firmware from \code{ti-linux-firmware}. Since this is outside
    the scope of Buildroot, we provide pre-built \code{tiboot3.bin}
    and \code{tispl.bin} in the lab data directory at
    \code{$HOME/__SESSION_NAME__-labs/buildroot-basic/data/}.

    If you want to learn how to build them yourself, refer to the
    Bootlin {\em Embedded Linux} training course, which covers the
    full AM62x boot chain in detail.

\end{itemize}

You're now done with the configuration!

\section{Building}

You could simply run \code{make}, but since we would like to keep a
log of the build, we'll redirect both the standard and error outputs
to a file, as well as the terminal by using the \code{tee} command:

\begin{bashinput}
make 2>&1 | tee build.log
\end{bashinput}

While the build is on-going, please go through the following sections
to prepare what will be needed to test the build results.

\section{Prepare the BeaglePlay}

\input{beagleplay-prepare.tex}

\section{Prepare a bootable micro-SD card}

The AM62x ROM code will look for boot images in a FAT partition on the
SD card. To be recognized by the ROM code, this partition must have a
special type and the {\em Bootable} flag set.

Let's prepare an SD card with such a partition.

Plug the SD card your instructor gave you on your workstation. Type
the \code{sudo dmesg} command to see which device is used by your
workstation. In case the device is \code{/dev/mmcblk0}, you will see
something like

\begin{verbatim}
[46939.425299] mmc0: new high speed SDHC card at address 0007
[46939.427947] mmcblk0: mmc0:0007 SD16G 14.5 GiB
\end{verbatim}

The device file name may be different (such as \code{/dev/sdb})
if the card reader is connected to a USB bus (either internally
or using a USB card reader).

In the following instructions, we will assume that your SD card is
seen as \code{/dev/mmcblk0} by your PC workstation.

Type the \code{mount} command to check your currently mounted
partitions. If SD partitions are mounted, unmount them:

\bashcmd{$ sudo umount /dev/mmcblk0p*}

We will erase the existing partition table by simply zero-ing the
first 16 MiB of the SD card:

\bashcmd{$ sudo dd if=/dev/zero of=/dev/mmcblk0 bs=1M count=16}

Now, let's use \code{cfdisk} to create the partitions that we need:

\begin{bashinput}
sudo apt install fdisk
sudo cfdisk /dev/mmcblk0
\end{bashinput}

If \code{cfdisk} asks you to \code{Select a label type}, choose
\code{dos}. This corresponds to traditional partition tables that
DOS/Windows would understand. \code{gpt} partition tables are needed
for disks bigger than 2 TB.

In the \code{cfdisk} interface, delete existing partitions, then
create two partitions with the following properties:

\begin{itemize}
\item First Partition
\begin{itemize}
  \item Size: \code{128MB}
  \item Select \code{primary} partition
  \item Type: \code{W95 FAT32 (LBA)} (\code{c} choice)
  \item Bootable flag enabled
\end{itemize}

\item Second Partition
\begin{itemize}
  \item Size: use the rest of the available space
  \item Select \code{primary} partition
  \item Type: \code{Linux} (\code{83} choice)
\end{itemize}
\end{itemize}

Press \code{Write} when you are done.

To make sure that partition definitions are reloaded on your
workstation, remove the SD card and insert it again.

Now create a FAT32 filesystem on the first partition and an ext4
filesystem on the second one:

\begin{bashinput}
sudo mkfs.vfat -F 32 -n boot /dev/mmcblk0p1
sudo mkfs.ext4 -L rootfs -E nodiscard /dev/mmcblk0p2
\end{bashinput}

You can now make your workstation automatically mount these partitions
by removing the SD card and plugging it back. They should now be
mounted on \code{/media/$USER/boot} and \code{/media/$USER/rootfs}.


Now everything should be ready. Hopefully by that time the Buildroot
build should have completed. If not, wait a little bit more.

\section{Flash the system}

Once Buildroot has finished building the system, it's time to put it
on the SD card:

\begin{itemize}

\item Copy the pre-built \code{tiboot3.bin} and \code{tispl.bin} from
  the lab data directory
  \code{$HOME/__SESSION_NAME__-labs/buildroot-basic/data/} to the
  \code{boot} partition of the SD card.

\item Copy \code{u-boot.img} from \code{output/images/} to the
  \code{boot} partition of the SD card.

\item Copy the \code{Image} and \code{k3-am625-beagleplay.dtb} files
  from \code{output/images/} to the \code{boot} partition of the SD
  card.

\item Extract the \code{rootfs.tar} file to the \code{rootfs}
  partition of the SD card, using:\\
  \inlinebash{sudo tar -C /media/$USER/rootfs/ -xf output/images/rootfs.tar}

\item Create a file named \code{extlinux/extlinux.conf} in the
  \code{boot} partition. This file should contain the following lines:

{\small
\begin{fileinput}
label buildroot
  kernel /Image
  devicetree /k3-am625-beagleplay.dtb
  append console=ttyS2,115200n8 root=/dev/mmcblk1p2 rootwait
\end{fileinput}
}

These lines teach the U-Boot bootloader how to load the Linux kernel
image and the Device Tree, before booting the kernel. It uses a
standard U-Boot mechanism called {\em distro boot command}, see
\url{https://source.denx.de/u-boot/u-boot/-/raw/master/doc/README.distro}
for more details.

\end{itemize}

Cleanly unmount the two SD card partitions, and eject the SD card.

\section{Boot the system}

Insert the SD card in the BeaglePlay's microSD slot. Press and hold
the \code{USR} button (located on the side of the board, near the
microSD slot) and power up the board through its USB-C connection.
You can release the \code{USR} button once you see output on the
serial console. Pressing \code{USR} forces the BeaglePlay to boot
from the SD card instead of from the internal eMMC.

You should see your system booting. Make sure that the U-Boot version
and build dates match with the current date. Do the same check for the
Linux kernel.

Login as \code{root} on the BeaglePlay, and explore the system. Run
\code{ps} to see which processes are running, and look at what
Buildroot has generated in \code{/bin}, \code{/lib}, \code{/usr} and
\code{/etc}.

Note: if your system doesn't boot as expected, make sure to reset the
U-Boot environment by running the following U-Boot commands:

\begin{bashinput}
env default -f -a
saveenv
\end{bashinput}

and reset. This is needed because the U-Boot loaded from the SD card
may still load the U-Boot environment from the eMMC. Ask your
instructor for additional clarifications if needed.

\section{Explore the build log}

Back to your build machine, since we redirected the build output to a
file called \code{build.log}, we can now have a look at it to see what
happened. Since the Buildroot build is quite verbose, Buildroot prints
before each important step a message prefixed by the \code{>>>}
sign. So to get an overall idea of what the build did, you can run:

\begin{bashinput}
grep ">>>" build.log
\end{bashinput}

You see the different packages between downloaded, extracted, patched,
configured, built and installed.

Feel free to explore the \code{output/} directory as well.
